% !PRINTED 712
\weberpage{712}
%<<CONT
\weq{2}{g_1 h_1 = \mathop{\mathrm{Lim}}\limits_{s=1} (s-1)\,\Phi_1(s),}
\weq{3}{\Phi_1(s) = \frac{1}{1-q^{-s}}\;\prod^{p_1}\frac{1}{(1-p_1^{-1/2\,sf})^{e}}\;\prod^{p_2}\frac{1}{(1-p_2^{-sf})^{1/2\,e}},}
in which the two product symbols refer to all the primes $p_1$ and
$p_2$. With the use of the symbols $f_1$, $e_1$ explained above one
can also write for this:
\weq{4}{\Phi_1(s) = \frac{1}{1-q^{-s}}\;\prod^{p}\frac{1}{(1-p^{-sf_1})^{e_1}},}
and extend the product symbol over \gesperrt{all} primes $p$.

The meaning of $g_1$ follows from \S.\,190, (11).

Now from the group $C$ of the characters $\chi$ we will single out a
divisor $C_1$, which shall consist of all those characters $\chi_1$
for which
\weq{5}{\chi_1(-1) = +1}
holds. It is clear at the outset that these characters $\chi_1$ form
a group $C_1$. But this group is not identical with $C$, for there
certainly is a character $\chi_0$ for which $\chi_0(-1) = -1$
holds. To obtain such a one we need only put, in \S.\,193,
(6), a primitive $\mu$th root of unity for $\Theta$,
or set $\varepsilon = -1$ in \S.\,193, (8). Then the
totality of the characters $\chi_0\chi_1$, none of which belong to
$C_1$, forms a subsidiary group $C_2$, whose elements may be
denoted by $\chi_2$. Every character $\chi$ is then
contained either in $C_1$ or in $C_2$; for if $\chi$ does not
occur in $C_1$, then $\chi_0^{-1}\chi$ occurs in it, and
consequently $\chi$ in $C_2$. Therefore $C_1$ is a divisor of $C$
of index 2.

The characters $\chi_1$ can be formed in the following way:

1) If $m$ is odd, we set, as in \S.\,193,
\[ n \equiv c^{\gamma} \pmod{m}, \]
and in
\weq{6}{\chi_1(n) = \Theta^{\gamma}}
we let $\Theta$ run through the totality of the $\tfrac12\mu$th
roots of unity.

2) If $m$ is even, we set
\[ n \equiv (-1)^{\alpha}5^{\beta}, \]
and obtain
\weq{7}{\chi_1(n) = \Theta^{\beta},}
in which $\Theta$ again runs through all the $\tfrac12\mu$th
roots of unity.
